\chapter{Legal, Ethical, \& Societal Implications}

\section{Legal Issues}

If MERIT is deployed, it must comply with data protection and privacy laws due to the processing of personal data such as CVs, LinkedIn profiles, and GitHub accounts. 
In particular, the system must comply with GDPR, meaning explicit consent will be required from candidates. 
Data must be stored securely, with mechanisms for data removal and correction. 
Furthermore, any automated decision-making processes must provide transparency and justification to explain how candidate scores were produced, reducing the risk of legal challenge.

\section{Ethical Issues}

MERIT must ensure fairness and transparency in candidate evaluation. 
There is a risk that the system may reinforce existing social biases present in historical data, especially if optional bias-mitigation extensions are not implemented. 
The system must respect candidate autonomy and avoid penalising applicants for missing or incomplete data, while providing clear explanations for screening outcomes to maintain trust and fairness.

\section{Professional Issues}

The use of MERIT raises considerations for HR practitioners and hiring managers. 
Users have a responsibility to understand the system's capabilities and limitations to avoid over-reliance on automated scores without appropriate human judgement. 
Misuse or misunderstanding of the system may result in poor hiring outcomes and presents reputational risk for organisations adopting the tool.

\section{Societal Issues}

MERIT may influence broader employment practices and labour market dynamics. 
As an automated screening system, it has the potential to reduce human bias and improve efficiency; however, it may also reinforce inequality if not carefully designed, particularly if certain groups are underrepresented within the data sources used. 
Societal trust in automated screening tools depends on transparency, accountability, and demonstrable fairness. 
Additionally, the system may impact perceptions of meritocracy, candidate experience, and organisational reputation.

